% 02-related.tex — 相关工作

\section{相关工作}
\label{sec:related}

\subsection{元胞自动机与迷宫}

Conway 1970 年提出 ``生命游戏'' 后, CA 在物理、生物学、计算机科学领域被广泛研究 \cite{conway1970}. Wolfram 2002 年的专著把 CA 动力学按行为分成 4 类 \cite{wolfram2002}, 其中 ``第 4 类'' (复杂类) 包含能产生持久结构与局部混沌的规则。生命游戏属于第 4 类, 但单条 B/S 字符串的搜索空间只有 $2^{18}$, 涌现迷宫结构的概率极低。

Mitchell, Hraber 与 Crutchfield 1993 年用遗传算法在 $2^{18}$ 的 B/S 字符串空间里搜索 ``密度稳态'' 规则 \cite{mitchell1993}, 发现了接近 0.5 密度的边缘动力学。Rejbrand 2017 年的网页笔记记录了用 8 邻域 + 优先级机制涌现迷宫的具体 B/S 组合 \cite{rejbrand2017}, 但仍是单族 B/S 字符串的扩展, 不支持族间仲裁。

McClendon 2001 年讨论了迷宫 ``难度'' 的形式化定义 (McClendon Difficulty), 用 $\gamma$ 度量死端墙密度 \cite{mcclendon2001}, 但与 Bellot 2021 的 $\nu$ 类似, 都是 ``越像迷宫值越大'' 的方向, 对几何反例同样失效。

\subsection{迷宫生成算法}

经典迷宫生成算法 (DFS, Kruskal, Prim, Growing Tree, Sidewinder, Binary Tree) 都能在 $O(n)$ 时间内生成 ``完美迷宫'' (无环, 连通, 任意两点恰好一条路径)。Pech 2002 年的硕士论文是早期综述 \cite{pech2002}, 详细比较了不同算法的走廊统计特征。

这些算法是 \textbf{真迷宫标签} 的来源 --- 本文 15-pattern 基准的 6 个 TRUE pattern 都由这些算法生成。

\subsection{几何度量}

Bellot 等人 2021 年的工作给出了 $F(M) = \nu(M) / \delta(M)$ 形式, 其中 $\nu$ 是 ``非显著墙'' (死端邻接类) 的计数, $\delta$ 是 ``twistiness proxy'' (diameter / diagonal) \cite{bellot2021}. 该论文用 9 个非迷宫反例测试 $F$, 报告 ``真迷宫值小, 伪迷宫值大'', 方向与直觉相反。本文 §4.4 在 15-pattern 基准上重复这个测试, 复现了 Bellot 论文报告的方向反转, 但同时报告 4/9 误判 (Spiral $F{=}11.85$ + 3 个 $F{=}0$ 图案) 排名低于真迷宫最低值。

神经架构搜索 (NAS) 中, DARTS \cite{liu2018darts} 与 Regularized Evolution \cite{real2019regularized} 都观察到 ``更大搜索空间不一定带来更好均值, 但峰值通常更高'' 的现象。本文 §6.1 在 (mask, maxFam) 消融中复现了同样的现象 (mf=1 均值 0.6984 高于 mf=8 均值 0.6306, 但 mf=8 拥有全场最高 0.8233)。

\subsection{WebGPU compute shader}

WebGPU 是 W3C 标准, 2023 年正式发布 compute shader 规范 \cite{webgpu2023}, 允许在浏览器内运行通用 GPU 计算。Skiena 2024 年讨论了 Codeforces 平台用 WebGPU 做大规模评估的实践 \cite{skiena2024}. 本文所有 sweep 的 CA 评估都在浏览器 WebGPU compute shader 上完成, 单 run $\sim$ 4 分钟 ($40 \times 60$) / $\sim$ 50 分钟 ($500 \times 2000$).

\subsection{本文与已有工作的关系}

本文与已有工作有 3 个明确的区别:
\begin{enumerate}[leftmargin=*, itemsep=0pt, topsep=0pt]
  \item \textbf{染色体结构。}已有工作集中在 B/S 字符串 ($2^{18}$ 候选), 本文 FamilyMask 把搜索空间扩展到 $2^{1648}$。
  \item \textbf{度量方向。}Bellot $F$ 是 ``值越小越像迷宫'' (方向反), 本文 $maze\_quality$ 是 ``值越大越像迷宫'' (方向正), 在 15-pattern 上达到 0/9 误判。
  \item \textbf{计算平台。}已有 CA 演化工作多在 CPU / 服务器 GPU 上跑, 本文在浏览器内通过 WebGPU compute shader 完成, 单机 0 部署成本。
\end{enumerate}
